% ID: SEL-D0026-OB-E04
% ESTADO: cotejada. Segunda pasada visual completa y borrador compilado revisados.
% Fuente: Selectividad/matematicas_aplicadas/suplemente_septiembre_EXAMEN 4-ANDALUCÍA.pdf
% SHA-256: c41ea8ffe46f2de3a4d7f82e43def97322e29653182ad3df54a57e2ead2beb19
% Página: 2; opción B; ejercicio original 4.

\textbf{(2.5 puntos)} La proporción de nacimientos que ocurren con luna llena en los hospitales de una ciudad se consideraba no inferior a 0.45, pero un estudio afirma que en la actualidad esta proporción ha descendido. Para contrastar esta hipótesis se han elegido al azar, en estos hospitales, a 200 recién nacidos, de los cuales 70 nacieron con luna llena. Decida mediante un contraste de hipótesis, con $H_0:\ p\geq 0.45$, si la afirmación del estudio es correcta con un nivel de significación del 1\%, indicando la región de rechazo.
